\documentclass[11pt]{article}  

% A template for problems and solution sets, by G. K. Vallis
% If possible, leave everything between the two dashed lines alone

% ----------------------------------------------------------------------


% The following two lines set the font. They are public domain
% and should be supported on all current TeX installations. 
% They may be commented out if they give you trouble. 
\usepackage[T1]{fontenc}
\usepackage{mathpazo}


% The following are some latex packages that you may find useful. All of
% them should be included in standard TeX installations. 
\usepackage{graphicx}
\usepackage{xspace}
\usepackage{bm}
\usepackage{paralist}
\usepackage[dvips]{geometry}
\usepackage{amsmath}
\usepackage{amssymb}


% Here are a few math abbreviations that may be useful. Note that math 
% variables should generally be set in italic, vectors in bold italic, and
% constants, such as e and i should be upright. Operators like the 'd' in 
% dx  and 'D' in D/Dt should also be upright. 
\newcommand{\ir}{\mathrm{i}}
\newcommand{\dr}{\mathrm{d}}
\newcommand{\D}{\mathrm{D}}
\newcommand{\e}{\, \mathrm{e}} \newcommand{\er}{\mathrm{e}}
\newcommand{\vb}{\bm {v}\xspace}
\newcommand{\ub}{\bm {u}\xspace}

% The following makes ODEs and PDEs easier to write.
% For an example, see the second problem below.
\newcommand{\dd}[3][]{{\dr^{#1} #2 \over \dr #3^{#1}}}
\newcommand{\pp}[3][]{{\partial^{#1} #2 \over \partial #3^{#1}}}
\newcommand{\DD}[1]{{\D#1 \over \D t}}

% Finally, the document geometry:
\geometry{centering, text={14cm,22cm}}
\renewcommand{\baselinestretch}{1.1}
\parindent 0pt

% ----------------------------------------------------------------------

% Start your own edits here.
% User-defined commands go here:



% Document text begins here:

\begin{document}
\centerline{\textbf{\textsf{\Large Solutions to problems in Chapter 5}}}
\medskip
\centerline{\textsf{\large by Your Name Goes Here} }
\medskip
\centerline{\today}

\vskip 1cm

This template may be used for either submitted problems or solutions. 
Any introductory text goes here. Note that each problem or solution is an
item in an enumerated list, and that the problem number goes in the square
brackets immediately following the word `item'. 

To typeset the document you may use either `pdflatex', which is the easiest
way, or `latex' followed by 'dvipdf' or `dvipdfm'. Please submit both the
latex source and the final pdf version of your solution. Thanks very much.

\begin{enumerate}
  
\item [5.6]  This is the (non)solution to problem 5.6. It contains some
text and some mathematics:
\begin{equation}
  \er^{\ir \pi} + 1 = 0, \qquad \psi = \widetilde{\psi} \e^{\ir(k x - \omega t)} .
\end{equation}


\item [5.9]  This is the (non)solution to problem 5.9, illustrating some
math constructs like pp and dd
\begin{equation}
 \dd \psi x + \pp[2]y x  = \int \chi \, \dr x ,
\end{equation}
and another, illustrating bold italic vectors:
\begin{equation}
 \DD \ub = \pp \ub t + \vb \cdot \nabla \ub .
\end{equation}

\item [5.10] Here is a problem or solution in two parts:
    \begin{enumerate}[(a)]
       \item Solution to part a. 
   
       \item Solution to part b
  \end{enumerate}


\end{enumerate}

Any concluding comments go here. 


\end{document}
